\begin{tablalafrox}{Diccionario de la cuenta 7.3. Fuente: (Bases Administrativas, arts. 17.2–17.6 y Formulario T-14, p. 63); (Bases Técnicas Transversales, secciones 17.2–17.6, pp. 32–34).}{tab:T14-dic-7-3}
  {L{1.2cm} Y Y Y C{0.9cm} L{2.3cm}}
  {\cab{Código} & \cab{Paquete de trabajo} & \cab{Entregable} & \cab{Criterio de aceptación} & \cab{Resp.} & \cab{Período}}
  7.3.1 & Innovación 3: documento de familias priorizadas y parámetros aprobados (INN-03.P1) & Documento de familias priorizadas y parámetros aprobados. & Los parámetros están aprobados por el Líder de Calidad y Pruebas y Calidad del CLIENTE. & DAT & E1 · Meses 3 a 5 \\
  7.3.2 & Innovación 3: asociación lote–sensor en el borde y en el camión (INN-03.P2) & Asociación lote–sensor en el borde y en el camión. & Un lote recorre todo su camino con historia completa en QA. & DAT & E1 · Meses 5 a 9 \\
  7.3.3 & Innovación 3: vista en M9 y orden de salida sugerido (INN-03.P3) & Vista en M9 y orden de salida sugerido. & No cambia la regla de Calidad y cierra antes del límite contractual del H4 (30-11-2027), sin integrar su entrega ni su acta. & DAT & E1 · Meses 9 y 10 (límite del H4) \\
  7.3.4 & Innovación 3: informe de validación en marcha blanca con Calidad (INN-03.P4) & Informe de validación en marcha blanca con Calidad. & Al mes 15 se cumplen dos metas: al menos el 95 \% de los lotes tiene historia completa, y al menos el 80 \% de las estimaciones queda dentro de ±15 \%. & DAT & E1 · Meses 11 a 15 (H5 y marcha blanca) \\
  7.3.5 & Innovación 3: revisión periódica de parámetros durante la Operación (INN-03.P5) & Parámetros ajustados con aprobación del Líder de Calidad de LafroX y de Calidad del CLIENTE. & Cada revisión periódica queda documentada y aprobada por el Líder de Calidad de LafroX y Calidad del CLIENTE. Si la exactitud baja del 80 \% dos veces seguidas, se hace un ajuste extraordinario. & SRE & Operación · Meses 24, 27, 30, 36, 39, 42, 48, 51 y 54 del proyecto. \\
\end{tablalafrox}
